\section{COAT: eliciting measured factors without assuming causal truth}
\label{exec19:coat-scope}

Consider a servicing team with customer narratives, offer decisions and later
repayment outcomes. Its first difficulty may be defining usable variables:
does a phrase describe a pre-existing cash-flow shock, a customer's reaction to
the offer, or a retrospective explanation written after default? A language
model can suggest a factor and extract its value, but the suggestion, the
measurement and the causal role require different evidence. The selected
13 October 2025 revision of Liu and colleagues' COAT framework addresses this
interface between unstructured observations and causal analysis. It proposes
and parses factors, then uses conventional causal-discovery feedback to refine
them. The paper separates COAT-MB for Markov blankets, COAT-PAG for refining
partial ancestral graphs, and COAT-ADJ for adjustment questions. This is a
useful architecture for a banking prototype whose variables are not already
settled.\citep{exec19_coat}

\subsection{What the representation assistant contributes}
\label{exec19:coat-representation}
In the paper's notation, $X$ is unstructured input, $\mathcal U$ is a set of
structured target variables and $W=h(X)$ is a proposed structured
representation. A factor instruction specifies a deterministic map from the
input to a numerical value. The initial sampling setup uses independently
drawn paired observations. COAT-MB proposes factors in rounds, parses their
values and updates a target's Markov blanket. Its feedback selects a cluster
where the existing factors leave high conditional entropy of the target. The
algorithm stops at a maximum number of rounds or an unchanged blanket.
These choices describe an adaptive search procedure; the stopping condition
does not itself certify that all relevant factors have been found.
\citep{exec19_coat}

For the bank, a factor definition therefore needs an observable interpretation,
an extraction rule, a timestamp and a missingness convention. ``Temporary
income interruption'' should not acquire a different meaning depending on the
desired intervention answer. If it is read from a note written after repayment,
its predictive usefulness does not make it a pre-offer confounder. Proposed
factors can be retained as hypotheses while their provenance and clock are
examined. An ordinary prediction improvement is useful feedback, but it does
not decide which variables are eligible adjustment controls.

The following project derivation makes an information limit explicit.
For a fixed mapping $W=h(X)$, if a
relevant variable $U$ is independent of $X$, it remains independent of $W$:
for any event $B$ of the representation,
\begin{equation}
 P(U=u,W\in B)
 =\sum_{x:h(x)\in B}P(U=u)P(X=x)
 =P(U=u)P(W\in B)
 \label{exec19:coat-information}
\end{equation}
in the discrete case. Naming a plausible hidden cause cannot create its
customer-specific measurement. Additional records or informative measurement
assumptions may help; the example is not a prohibition on extracting genuine
information already present in text. The paper's ``hidden world'' concerns
eliciting high-level factors from observations, not an unconditional ability
to identify every unobserved cause.\citep{exec19_coat}

\subsection{The parsing assumption is scientifically substantive}
\label{exec19:coat-parsing}
The selected paper explicitly assumes faithful and Markov correspondence
between a graph and the factor distribution, with conditional independences
preserved after parsing. It acknowledges that annotation error can violate
this premise. High marginal accuracy is consequently insufficient to establish
that a parsed variable can replace its true counterpart in adjustment or
conditional-independence tests.\citep{exec19_coat}

An exact project-authored mechanism illustrates the issue. Let $U$ be a fair
binary pre-offer cause, let the offer $A$ and repayment $Y$ be conditionally
independent given $U$, and specify
\begin{equation}
 P(A=1\mid U)=P(Y=1\mid U)=\tfrac14+\tfrac12U,
 \qquad W=U\mathbin{\oplus}E,
 \qquad E\sim\operatorname{Bernoulli}(1/4),
 \label{exec19:coat-noise-law}
\end{equation}
where $E$ is independent of the other structural randomizations and $\oplus$
denotes exclusive-or. There is no $A\to Y$ arrow: the true offer effect is
zero, and both interventional repayment risks equal $1/2$. The proxy $W$ is
75\% accurate. Nevertheless, conditional on either value of $W$, $U$ retains
variance $3/16$. The conditional covariance is
\begin{equation}
 \operatorname{Cov}(A,Y\mid W=w)
 =\tfrac14\operatorname{Var}(U\mid W=w)=\tfrac3{64}>0.
 \label{exec19:coat-covariance}
\end{equation}
The factor error has destroyed the conditional independence that held given
$U$. This is a deliberate violation of the paper's parsing assumption, not
evidence of failure when that assumption holds.\citep{exec19_coat_reference}

The exact conditional repayment risks for $(a,w)=(0,0),(1,0),(0,1),(1,1)$
are $3/10,1/2,1/2,7/10$. Standardizing over the fair proxy gives
\begin{equation}
 \sum_w P(Y=1\mid A=0,w)P(w)=\tfrac25,
 \qquad
 \sum_w P(Y=1\mid A=1,w)P(w)=\tfrac35.
 \label{exec19:coat-proxy-effect}
\end{equation}
The spurious effect is $1/5$. With an assumed USD 100 payoff per incremental
repayment and USD 10 offer cost, that calculation yields USD 10, whereas the
true zero-effect offer has net value USD $-10$. The positive joint law is
enumerated with rational arithmetic; no COAT run, empirical annotation rate
or actual bank benefit is being reported. The proposed implication is to
evaluate extraction errors by variable, outcome, treatment and time, rather
than accepting a single overall accuracy number as a causal qualification.
\citep{exec19_coat_reference}

\subsection{Separate iterative discovery from a finite stopping guarantee}
\label{exec19:coat-attempts}
COAT's ability definition posits a positive probability of proposing a useful
factor and a positive multiplicative reduction in residual information. Its
convergence discussion then uses these quantities to characterize successful
discovery rounds. The proof of Proposition 5 treats the number of useful rounds
as binomial and explicitly approximates that distribution by a Gaussian.
That calculation is informative about the proposed regime, but an
uncontrolled approximation cannot be adopted as an exact finite probability
certificate for a banking run. The model for repeated useful proposals also
needs to match an adaptive agent's history and failure modes.\citep{exec19_coat}

For example, four independent proposal rounds with success probability $1/2$
have no-success probability $1/16$. If all four outcomes equal one shared
fair Bernoulli variable, each round still has marginal success probability
$1/2$, but the chance of no success is $1/2$. Repeated versions of the same
unsuccessful hypothesis can resemble the latter dependence pattern. The
authored example does not diagnose a measured COAT dependence structure; it
shows why marginal responsiveness alone cannot specify a multi-round success
probability.\citep{exec19_coat_reference}

A public prototype should predeclare the proposal budget, the meaning of a
useful factor and the evidence supplied as feedback. If a quantitative
convergence claim is sought, it must state whether successes are independent
with a fixed probability or satisfy an explicit conditional lower bound given
the preceding history. It also needs a justified contraction statement and an
exact or controlled probability bound. Without those premises, a maximum-round
stop is a resource limit and an unchanged blanket is a search observation.
Neither is upgraded to ``all causal variables recovered.'' Separate final
tasks should remain outside the feedback loop so that repeated improvement on
the development sample is not counted again as fresh validation.

\subsection{Graph preservation and effect estimation retain their conditions}
\label{exec19:coat-graph}
The paper's graph extensions add further structure beyond a single target's
blanket. Its Proposition 8 states adjacency and arrowhead preservation within
a subset under blanket inclusion in the enlarged variable set and the stated
maximally informative PAG setup. The adjustment development additionally uses
amenability, forbidden-variable and blocking conditions. Proposition 15
preserves an adjustment set under its specified blanket-inclusion premise.
These are conditional results about the constructed measured-variable graph;
they are not general certificates for arbitrary hidden-confounder or
cross-world questions. A proposed adapter must preserve uncertain marks and
refuse an effect query when the applicable premises or adjustment route are
unresolved.\citep{exec19_coat}

The marketing evaluation uses Hillstrom and Retail Hero in a confounded
evaluation setting with synthesized social-media text. For each experiment,
the authors choose the largest valid adjustment sets from the generated PAG
and average estimated effects when there are multiple such sets; experiments
are repeated three times. The extracted Table 7 reports COAT RMSEs of
0.47 and 0.50 for DeepSeek-V3, and 1.90 and 0.78 for GPT-4.1-mini on the
two datasets. The human-factor N-IPW reference reports 1.32 and 1.39.
Thus the comparative result depends on model and comparator: the first
GPT-4.1-mini number does not beat that human-factor reference. These are
source-reported outcomes in the paper's design, not locally reproduced
estimation or confidence-interval calibration.\citep{exec19_coat}

For a proposed lending study, average error, graph validity, adjustment
availability and interval coverage should be recorded separately. A failure
to recover a valid set must remain in the request denominator. Averaging
effects over several selected sets does not remove the need to justify their
validity or account for selection and dependence. A generated-text benchmark
can expose parser and graph failures at low cost, but it does not establish
that authentic customer narratives measure the same factors.

\subsection{Financial examples motivate a clock and cost audit}
\label{exec19:coat-finance}
COAT's financial case uses Microsoft price and news observations from
2006--2009, with the first 200 trading days for elicitation and the next 400
for evaluation. The target describes the sign of a following-week return.
The reported strategy fits a rolling regression and takes a long or short
position according to its predicted sign. This is an observational trading
illustration. Its reported performance does not identify the effect of a bank
action on a customer outcome. The printed rolling-window summation limits
also need clarification against the executable source; this proposal does not
silently repair them and assume the code did the same thing.\citep{exec19_coat}

A faithful public reproduction would bind news-publication times, data
availability, revision policy, normalization windows, forecast horizon and
trade execution time. Overlapping weekly outcomes require a declared
dependence treatment. It would retain transaction costs, financing and
failed-query costs before interpreting a trading return as a useful economic
comparison. This is an inexpensive way to test the proposed workflow's
meaning before requesting bank data, rather than treating a published
historical return as bank ROI.

The proposed integration is therefore a staged representation-assistant
boundary: factor instructions and source passages enter a reviewed ontology;
the parser produces values with provenance and explicit missingness; graphical
analysis preserves uncertainty and scientific refusals; downstream estimation
records its own assumptions. The paper's project link is a source-code lead,
not an audited repository, license or runtime. Native artifact review,
measurement-error evaluation and finite stopping guarantees remain separate
work. Public or stipulated gold factors can demonstrate whether this route
merits those later experiments, while independent acceptance and authentic
bank truth remain later evidence.\citep{exec19_coat,exec19_coat_reference}
